\documentclass[12pt, titlepage]{article}

\usepackage{booktabs}
\usepackage{tabularx}
\usepackage{hyperref}
\hypersetup{
    colorlinks,
    citecolor=blue,
    filecolor=black,
    linkcolor=red,
    urlcolor=blue
}
\usepackage[round]{natbib}

\input{../Comments.text}
\input{../Common.text}

\begin{document}

\title{System Verification and Validation Plan for \progname{}} 
\author{\authname}
\date{\today}
	
\maketitle

\pagenumbering{roman}

\section*{Revision History}

\begin{tabularx}{\textwidth}{p{3cm}p{2cm}X}
\toprule {\bf Date} & {\bf Version} & {\bf Notes}\\
\midrule
Date 1 & 1.0 & Notes\\
Date 2 & 1.1 & Notes\\
\bottomrule
\end{tabularx}

~\\
\wss{The intention of the VnV plan is to increase confidence in the software.
However, this does not mean listing every verification and validation technique
that has ever been devised.  The VnV plan should also be a \textbf{feasible}
plan. Execution of the plan should be possible with the time and team available.
If the full plan cannot be completed during the time available, it can either be
modified to ``fake it'', or a better solution is to add a section describing
what work has been completed and what work is still planned for the future.}

\wss{The VnV plan is typically started after the requirements stage, but before
the design stage.  This means that the sections related to unit testing cannot
initially be completed.  The sections will be filled in after the design stage
is complete.  the final version of the VnV plan should have all sections filled
in.}

\newpage

\tableofcontents

\listoftables
\wss{Remove this section if it isn't needed}

\listoffigures
\wss{Remove this section if it isn't needed}

\newpage

\section{Symbols, Abbreviations, and Acronyms}

\renewcommand{\arraystretch}{1.2}
\begin{tabular}{l l} 
  \toprule		
  \textbf{symbol} & \textbf{description}\\
  \midrule 
  T & Test\\
  \bottomrule
\end{tabular}\\

\wss{symbols, abbreviations, or acronyms --- you can simply reference the SRS
  \citep{SRS} tables, if appropriate}

\wss{Remove this section if it isn't needed}

\newpage

\pagenumbering{arabic}

This document ... \wss{provide an introductory blurb and roadmap of the
  Verification and Validation plan}

\section{General Information}

\subsection{Summary}

\wss{Say what software is being tested.  Give its name and a brief overview of
  its general functions.}

\subsection{Objectives}

\wss{State what is intended to be accomplished.  The objective will be around
  the qualities that are most important for your project.  You might have
  something like: ``build confidence in the software correctness,''
  ``demonstrate adequate usability.'' etc.  You won't list all of the qualities,
  just those that are most important.}

\wss{You should also list the objectives that are out of scope.  You don't have 
the resources to do everything, so what will you be leaving out.  For instance, 
if you are not going to verify the quality of usability, state this.  It is also 
worthwhile to justify why the objectives are left out.}

\wss{The objectives are important because they highlight that you are aware of 
limitations in your resources for verification and validation.  You can't do everything, 
so what are you going to prioritize?  As an example, if your system depends on an 
external library, you can explicitly state that you will assume that external library 
has already been verified by its implementation team.}

\subsection{Custom Documents}

\wss{Summarize the CDs (Custom Documents) that were tackled by this project.  CDs 
can include usability testing, code walkthroughs, user documentation, formal
proof, GenderMag personas, Design Thinking, etc.  CDs should have already
been approved by the course instructor as included in your problem statement.
You can use a pull request to update your CDs (in TeamComposition.csv or
Repos.csv) if your plan changes as a result of the VnV planning exercise.}

\subsection{Relevant Documentation}

\wss{Reference relevant documentation.  This will definitely include your SRS
  and your other project documents (design documents, like MG, MIS, etc).  You
  can include these even before they are written, since by the time the project
  is done, they will be written.  You can create BibTeX entries for your
  documents and within those entries include a hyperlink to the documents.}

\citet{SRS}

\wss{Don't just list the other documents.  You should explain why they are relevant and 
how they relate to your VnV efforts.}

\section{Plan}

\wss{Introduce this section.  You can provide a roadmap of the sections to
  come.}

\subsection{Verification and Validation Team}

\wss{Your teammates.  Maybe your supervisor.
  You should do more than list names.  You should say what each person's role is
  for the project's verification.  A table is a good way to summarize this information.}

\subsection{SRS Verification}

\wss{List any approaches you intend to use for SRS verification.  This may
  include ad hoc feedback from reviewers, like your classmates (like your
  primary reviewer), or you may plan for something more rigorous/systematic.}

\wss{If you have a supervisor for the project, you shouldn't just say they will
read over the SRS.  You should explain your structured approach to the review.
Will you have a meeting?  What will you present?  What questions will you ask?
Will you give them instructions for a task-based inspection?  Will you use your
issue tracker?}

\wss{Maybe create an SRS checklist?}

\subsection{Design Verification}

\wss{Plans for design verification}

\wss{The review will include reviews by your classmates}

\wss{Create a checklists?}

\subsection{Verification and Validation Plan Verification}

\wss{The verification and validation plan is an artifact that should also be
verified.  Techniques for this include review and mutation testing.}

\wss{The review will include reviews by your classmates}

\wss{Create a checklists?}

\subsection{Implementation Verification}

\wss{You should at least point to the tests listed in this document and the unit
  testing plan.}

\wss{In this section you would also give any details of any plans for static
  verification of the implementation.  Potential techniques include code
  walkthroughs, code inspection, static analyzers, etc.}

\wss{The final class presentation in CAS 741 could be used as a code
walkthrough.  There is also a possibility of using the final presentation (in
CAS741) for a partial usability survey.}

\subsection{Automated Testing and Verification Tools}

\wss{What tools are you using for automated testing.  Likely a unit testing
  framework and maybe a profiling tool, like ValGrind.  Other possible tools
  include a static analyzer, make, continuous integration tools, test coverage
  tools, etc.  Explain your plans for summarizing code coverage metrics.
  Linters are another important class of tools.  For the programming language
  you select, you should look at the available linters.  There may also be tools
  that verify that coding standards have been respected, like flake9 for
  Python.}

\wss{If you have already done this in the development plan, you can point to
that document.}

\wss{The details of this section will likely evolve as you get closer to the
  implementation.}

\subsection{Software Validation}

\wss{If there is any external data that can be used for validation, you should
  point to it here.  If there are no plans for validation, you should state that
  here.}

\wss{You might want to use review sessions with the stakeholder to check that
the requirements document captures the right requirements.  Maybe task based
inspection?}

\wss{For those capstone teams with an external supervisor, the Rev 0 demo should 
be used as an opportunity to validate the requirements.  You should plan on 
demonstrating your project to your supervisor shortly after the scheduled Rev 0 demo.  
The feedback from your supervisor will be very useful for improving your project.}

\wss{For teams without an external supervisor, user testing can serve the same purpose 
as a Rev 0 demo for the supervisor.}

\wss{This section might reference back to the SRS verification section.}

\section{System Tests}

\wss{There should be text between all headings, even if it is just a roadmap of
the contents of the subsections.}

\subsection{Tests for Functional Requirements}

\wss{Subsets of the tests may be in related, so this section is divided into
  different areas.  If there are no identifiable subsets for the tests, this
  level of document structure can be removed.}

\wss{Include a blurb here to explain why the subsections below
  cover the requirements.  References to the SRS would be good here.}

\subsubsection{Area of Testing1}

\wss{It would be nice to have a blurb here to explain why the subsections below
  cover the requirements.  References to the SRS would be good here.  If a section
  covers tests for input constraints, you should reference the data constraints
  table in the SRS.}
		
\paragraph{Title for Test}

\begin{enumerate}

\item{test-id1\\}

Control: Manual versus Automatic
					
Initial State: 
					
Input: 
					
Output: \wss{The expected result for the given inputs.  Output is not how you
are going to return the results of the test.  The output is the expected
result.}

Test Case Derivation: \wss{Justify the expected value given in the Output field}
					
How test will be performed: 
					
\item{test-id2\\}

Control: Manual versus Automatic
					
Initial State: 
					
Input: 
					
Output: \wss{The expected result for the given inputs}

Test Case Derivation: \wss{Justify the expected value given in the Output field}

How test will be performed: 

\end{enumerate}

\subsubsection{Area of Testing2}

...

\subsection{Tests for Nonfunctional Requirements}

\wss{The nonfunctional requirements for accuracy will likely just reference the
  appropriate functional tests from above.  The test cases should mention
  reporting the relative error for these tests.  Not all projects will
  necessarily have nonfunctional requirements related to accuracy.}

\wss{For some nonfunctional tests, you won't be setting a target threshold for
passing the test, but rather describing the experiment you will do to measure
the quality for different inputs.  For instance, you could measure speed versus
the problem size.  The output of the test isn't pass/fail, but rather a summary
table or graph.}

\wss{Tests related to usability could include conducting a usability test and
  survey.  The survey will be in the Appendix.}

\wss{Static tests, review, inspections, and walkthroughs, will not follow the
format for the tests given below.}

\wss{If you introduce static tests in your plan, you need to provide details.
How will they be done?  In cases like code (or document) walkthroughs, who will
be involved? Be specific.}

\subsubsection{Area of Testing1}
		
\paragraph{Title for Test}

\begin{enumerate}

\item{test-id1\\}

Type: Functional, Dynamic, Manual, Static etc.
					
Initial State: 
					
Input/Condition: 
					
Output/Result: 
					
How test will be performed: 
					
\item{test-id2\\}

Type: Functional, Dynamic, Manual, Static etc.
					
Initial State: 
					
Input: 
					
Output: 
					
How test will be performed: 

\end{enumerate}

\subsubsection{Safety-Critical Requirements}
\label{sec:safety-tests}

These tests cover the 18 safety-critical requirements in SRS Section~12.2,
which come from the \href{https://github.com/AHMEDELZARIA/squigglr/blob/main/docs/HazardAnalysis/HazardAnalysis.tex}{Hazard Analysis}. A failure here could harm a
child's well-being, privacy or records, so each requirement gets two
different kinds of evidence: a test below, and the safety inspection
described next. Every test below is dynamic except T-SAF-004, which checks
wording in templates and so is an inspection; T-SAF-016 combines a code
review with recordings on the iPad. The tests try to break each safeguard, for example by
having a child tap adult-only controls, cutting the connection during an
attempt, or switching between children on a shared device. All tests use
test accounts and prepared data only; no child takes part (SRS CON-005).
In these tests, ``the adult'' means the parent or educator responsible for
the child, and ``the adult check'' means the check in \texttt{ADULT\_CHECK}
(for example a PIN or password) that only the adult can complete.

\paragraph{Safety inspection (SI).} Before each revision (Rev~0 and
Rev~1), the code that implements each requirement below is inspected by two
team members who did not write it. The Backend Lead (Rayan Nasrallah) runs
the inspection and the Testing Lead (Muhammad Huzaifah) records the results.
For each requirement, the inspectors:
\begin{enumerate}
\item find every code path that triggers the safeguard, using the
test-to-requirement traceability table;
\item check that no path skips it, for example an adult-only action reachable
from a URL or a save that reports success before the server responds;
\item check that a dynamic test exercises each path, and add one if not;
\item check the Code checklist items on error handling and input validation;
\item record each finding as a GitHub issue labelled \texttt{safety}, and
re-inspect after the fix.
\end{enumerate}
The inspection passes when every requirement has no open \texttt{safety}
issue. Results are reported in the VnV Report.

\paragraph{Why this is enough.} Testing cannot prove that a safeguard always
works; it can only raise confidence that it does. One test per
requirement, each with cases chosen to break the safeguard, plus an
inspection that looks for paths the test missed, gives two independent
checks. A second dynamic test is added wherever the inspection finds an
untested path. The inspection covers only the code tied to these 18
requirements, because a full review of every module by two people is not
realistic in the time the team has. The rest of the code is covered by the
code checklist and the normal PR review.

\paragraph{How a lesson works.} Several of these tests check the limits in a
guided lesson, so this paragraph summarizes the lesson flow from SRS BR-002,
UC-002 to UC-005 and the state machine in SRS Section~9.2. The adult picks the
child and the lesson and completes the adult check. The child then watches a
demonstration and works through three stages in order: tracing a full
outline, drawing between endpoints, and drawing with no guide. An attempt
that meets \texttt{STAGE\_PASS} moves the child to the next stage
automatically, and a pass at the last stage completes the lesson. Failed
attempts are counted \emph{per stage}, and the count starts again whenever
the stage changes:
\begin{itemize}
\item after \texttt{DIFFICULTY\_STREAK} failed attempts in a row, Squigglr
suggests to the adult that the child go back a stage; the stage changes only
if the adult completes the adult check and accepts, and the child keeps
practising while the suggestion waits;
\item after \texttt{ATTEMPT\_LIMIT} failed attempts at one stage, the lesson
ends automatically, unless the child is improving, in which case the child
gets up to \texttt{ATTEMPT\_EXTENSION} more attempts, never more than
\texttt{ATTEMPT\_LIMIT\_MAX} in total at that stage;
\item after \texttt{LESSON\_IDLE\_PROMPT} with no input, Squigglr shows the
child a check-in, and after \texttt{LESSON\_IDLE\_END} with no input, the
lesson ends; time without input never counts as a failed attempt.
\end{itemize}
Because \texttt{ATTEMPT\_LIMIT} is always larger than
\texttt{DIFFICULTY\_STREAK} (SAF-RAY-008), the adult always receives the
suggestion before the lesson ends: the suggestion is the early warning and
the limit is the backstop. The adult sets these values for each child within
ranges set by the team. The examples below use illustrative values
(\texttt{DIFFICULTY\_STREAK} = 3, \texttt{ATTEMPT\_LIMIT} = 6,
\texttt{ATTEMPT\_EXTENSION} = 2, \texttt{ATTEMPT\_LIMIT\_MAX} = 8) to show
the expected behaviour; the tests themselves use the ends of each range.

\begin{enumerate}

\item{T-SAF-001: Adult-only actions need the adult check\\}

Purpose: choosing a lesson, choosing a starting stage, switching child and
accepting a go-back suggestion belong to the adult, so a child must not be
able to complete any of them.

Type: Dynamic, manual and automatic. Covers SAF-RAY-001.

Initial State: an adult account with two child profiles; child A is in a
guided lesson, and a go-back suggestion is waiting for the adult.

Input: on the lesson screen, the tester (acting as the child) taps every
visible control, taps the screen 50 times at random positions, opens the
web address of each adult-only action directly, uses the browser's back
and forward buttons and reloads the page. Each time the adult check
appears, the tester leaves it or enters a wrong answer. Finally, the tester
(acting as the adult) completes the adult check and accepts the suggestion.

Output: every adult-only action asks for the adult check, and none completes
until the check is passed. The suggestion stays waiting and the child can
keep drawing at the current stage. After the adult passes the check and
accepts, the child moves back one stage.

Example: child A fails three attempts in a row at the endpoints stage, and a
go-back suggestion appears. Child A taps ``accept''; Squigglr asks for the
adult check, which child A cannot complete, so the child stays at the
endpoints stage and keeps practising. The adult then completes the check and
accepts, and child A moves back to the tracing stage.

Test Case Derivation: SAF-RAY-001 and its hazards H-RAY-001 and H-RAY-006.
Random taps imitate a child who cannot read exploring the screen.

How test will be performed: a Playwright script performs the taps, web
addresses and navigation and checks the lesson state after each one; a
tester repeats the taps by hand with the stylus on the iPad.

\item{T-SAF-002: Notice on an unsupported device, browser or input\\}

Purpose: drawing results may be unreliable on a device or browser that is not
in \texttt{SUPPORTED\_DEVICES}, or when the child draws with a finger instead
of a stylus, so the adult must be told and can decide whether to continue.

Type: Dynamic, automatic and manual. Covers SAF-RAY-002.

Initial State: the adult has chosen a lesson; it has not started.

Input: (a) the lesson is started once with each combination in
\texttt{SUPPORTED\_DEVICES}, and once with each of three combinations that
are not in it: an older browser version, a phone-sized screen, and a browser
without pointer event support; (b) on the iPad, a supported lesson is
started and the first stroke is drawn with a finger; (c) the same, with the
first stroke drawn with the stylus.

Output: (a) for each unsupported combination, a notice the adult can read
appears before the first attempt; for each supported combination, no notice
appears; (b) the notice appears after the first stroke; (c) no notice
appears.

Test Case Derivation: the fit criterion of SAF-RAY-002 (H-RAY-002). A
browser cannot tell whether a stylus is connected before it is used, so
case (a) checks only what the browser reports (version, screen size and
pointer event support), and cases (b) and (c) check the input type the
browser reports for the first stroke (pen or touch).

How test will be performed: for (a), Playwright sets the browser version,
screen size and pointer event support for each combination and checks for
the notice; for (b) and (c), a tester draws on the iPad, once with a finger
and once with the stylus.

\item{T-SAF-003: Attempts kept until the server confirms them\\}

Purpose: an attempt must never look saved when it is not, or the adult would
make decisions from incomplete records.

Type: Dynamic, automatic. Covers SAF-RAY-003 and SAF-SHA-003.

Initial State: a guided lesson is open and online.

Input: (a) the connection is cut during an attempt; (b) the connection is cut
after the attempt is finished but before the server confirms the save, and
the page is then reloaded, and separately closed and reopened, while
offline; (c) the server returns an error for \texttt{FAILED\_SAVES} saves in
a row; (d) \texttt{FAILED\_SAVES} saves succeed normally.

Output: in (a) to (c), no attempt is shown as saved before the server
confirms it, the adult is told which attempts are not saved, and every
attempt is saved once the connection or server is back; in (d), no failure
message appears.

Test Case Derivation: SAF-RAY-003 and SAF-SHA-003 (H-RAY-002, H-SHA-003).
Case (d) checks that the warning does not appear when nothing failed.

How test will be performed: Playwright turns the network off and makes the
server return errors, then compares the server's records with the attempts
made.

\item{T-SAF-004: No diagnosis or comparison in summaries and exports\\}

Purpose: Squigglr is a practice tool, not an assessment, so no summary or
export may label, diagnose or compare a child, and each must say so.

Type: Static (inspection), manual. Covers SAF-RAY-004.

Initial State: the templates for every summary and export at that revision.

Input: one team member who did not write the templates reads each one.

Output: every template states that it describes practice only and is not an
assessment or diagnosis, and none uses a diagnostic label, judging wording
(for example ``struggling'', ``behind'' or ``below average'') or a comparison
with other children.

Test Case Derivation: the fit criterion of SAF-RAY-004 (H-RAY-003), which
asks for an inspection of every summary and export template. The team does
not intend to add diagnoses, so the realistic risks are a new template that
leaves out the required statement and wording that judges the child; a read
of each template catches both.

How test will be performed: once per revision (Rev~0 and Rev~1), as part of
the safety inspection; any problem is recorded as a GitHub issue labelled
\texttt{safety}.

\item{T-SAF-005: Lesson ends automatically at the attempt limit\\}

Purpose: a child must not be left repeating a stage they keep failing, so
the lesson must end by itself at the limit the adult set.

Type: Dynamic, automatic. Covers SAF-RAY-005.

Initial State: a child profile with \texttt{ATTEMPT\_LIMIT} and
\texttt{ATTEMPT\_LIMIT\_MAX} at the lowest values in their ranges; the test
is then repeated with the highest values.

Input: (a) at the tracing stage, the failing attempt
\texttt{circle-\allowbreak{}fail.json} is replayed until the lesson ends; (b)
the child passes tracing after a few failures, and the failing attempt is
replayed at the endpoints stage until the lesson ends; (c) a sequence of
attempts that improves enough to earn an extension is replayed until the
lesson ends.

Output: the lesson ends after exactly \texttt{ATTEMPT\_LIMIT} failed
attempts at one stage, or \texttt{ATTEMPT\_LIMIT\_MAX} when extended, with no
action from the child; failures at an earlier stage do not count at the new
stage; and the saved record matches the attempts made, including the
furthest stage reached.

Example: the sixth failed tracing attempt ends the lesson. In case (b), a
child who passed tracing after four failures still gets six attempts at the
endpoints stage. In case (c), an improving child gets attempts 7 and 8, and
the lesson ends after the eighth failure.

Test Case Derivation: SAF-RAY-005 (H-RAY-004). Testing both ends of the
ranges follows the supervisor's advice to test at the minimum and maximum
(SRC-008).

How test will be performed: Playwright counts the attempts and checks that
no attempt is accepted after the lesson ends.

\item{T-SAF-006: Limit endings are shown to the adult across repeated lessons\\}

Purpose: each time a lesson is done, it is a new lesson record, so the adult
must see how often the lesson ended at the limit across all of them; doing
the lesson again must not hide earlier failures.

Type: Dynamic, automatic. Covers SAF-RAY-006.

Initial State: as in T-SAF-005.

Input: the adult starts the lesson and it is run until it ends at the
limit; the adult then starts the same lesson again, and this is repeated so
the lesson is done three times. In the third run, an improving sequence
earns an extension before the lesson ends at the limit.

Output: the adult's view of that lesson shows three lesson records, each
ending at the limit, one extension, and every attempt from all three
records.

Example: the adult sees ``ended at the limit: 3 times; extended: once'' and
all 6 + 6 + 8 = 20 attempts.

Test Case Derivation: SAF-RAY-006 (H-RAY-005).

How test will be performed: the script compares the counts and the list of
attempts in the adult's view with the runs made.

\item{T-SAF-007: Attempts saved only to the active child\\}

Purpose: on a shared device, one child's attempts must never be saved to
another child's record.

Type: Dynamic, automatic and manual. Covers SAF-RAY-007.

Initial State: one adult account with three child profiles (A, B and C) on
one device.

Input: the adult switches between the children in the order A, B, A, C,
and each child makes two attempts; the connection is cut once during child
B's turn.

Output: every practice screen shows the active child's avatar, and every
attempt and suggestion is saved under the child whose avatar was shown when
it was made.

Test Case Derivation: SAF-RAY-007 (H-RAY-006). The lost connection checks
that a delayed save still goes to the right child.

How test will be performed: Playwright checks each child's records; a
tester checks by hand that each avatar can be recognized without reading.

\item{T-SAF-008: Adult-set values stay within their ranges\\}

Purpose: a mistyped or extreme value could give a child an impossible
lesson, so each value must stay within the range the team set and only the
adult may change it.

Type: Dynamic, automatic. Covers SAF-RAY-008.

Initial State: a new child profile with no values chosen.

Input: (a) no value is set; (b) each value is set to its minimum, its
maximum, one below the minimum and one above the maximum; (c)
\texttt{ATTEMPT\_LIMIT} is set equal to \texttt{DIFFICULTY\_STREAK}, and
\texttt{ATTEMPT\_LIMIT\_MAX} is set one below \texttt{ATTEMPT\_LIMIT}; (d) a
change is tried without completing the adult check.

Output: (a) each value is the middle of its range; (b) in-range values are
saved and out-of-range values are refused; (c) both settings are refused;
(d) nothing changes. The values in use are shown with the child's records.

Example: if the team-set range for \texttt{ATTEMPT\_LIMIT} were 4 to 10, an
unset value would be 7, values 4 and 10 would be saved, and 3 and 11 would
be refused. With \texttt{DIFFICULTY\_STREAK} = 3, an \texttt{ATTEMPT\_LIMIT}
of 3 is refused, because the lesson would end before the adult received the
go-back suggestion.

Test Case Derivation: the fit criterion of SAF-RAY-008 (H-RAY-007), using
boundary values on each side of each range.

How test will be performed: a Vitest test checks the validation function
with every case, and Playwright repeats cases (b) to (d) through the
settings screen and through direct server requests.

\item{T-SAF-009: Duplicate input events counted once\\}

Purpose: a stylus event reported twice must not change a drawing's result,
or a child could pass or fail for no real reason.

Type: Dynamic, automatic. Covers SAF-AHM-001.

Initial State: a guided lesson at the tracing stage.

Input: \texttt{square-\allowbreak{}4-\allowbreak{}strokes.json} is replayed
as recorded, then with every input event sent twice; then
\texttt{overlap-\allowbreak{}2-\allowbreak{}strokes.json}, two separate
strokes that cross, is replayed.

Output: the recorded points, classification and pass result are the same
with and without the duplicates, and both crossing strokes are kept.

Test Case Derivation: the fit criterion of SAF-AHM-001 (H-AHM-002).

How test will be performed: Playwright replays both versions and compares
the saved records.

\item{T-SAF-010: Drawing kept when the canvas changes size\\}

Purpose: rotating the iPad or resizing the window must not distort a drawing
or change its result.

Type: Dynamic, automatic. Covers SAF-AHM-002.

Initial State: a guided lesson with an attempt partly drawn.

Input: the window is resized and the iPad is rotated during the attempt;
then input events with positions outside the canvas, missing positions
(\texttt{NaN}) and very large values are sent.

Output: the order and relative positions of the points are unchanged, and
the classification is the same as without resizing; none of the invalid
points is counted or makes the attempt pass.

Test Case Derivation: the fit criterion of SAF-AHM-002 (H-AHM-003).

How test will be performed: Playwright resizes the drawing area and compares
the scaled points with the original; a Vitest test sends the invalid values
to the input handler.

\item{T-SAF-011: Every data request checks the account\\}

Purpose: one family's or school's data must never be reachable from another
account.

Type: Dynamic, automatic. Covers SAF-SHA-001, with T-SEC-001.

Initial State: two adult accounts, A and B, each with a child profile and
saved attempts.

Input: logged in as A, the tester sends requests to read, change and delete
B's child profiles, attempts, settings and summaries to every server
endpoint, then repeats them with no session and with an expired session.

Output: every request is refused, and none of B's data is returned or
changed.

Test Case Derivation: the fit criterion of SAF-SHA-001 (H-SHA-001).

How test will be performed: a script lists every endpoint from the code and
sends each request; the list is checked in the safety inspection, so a new
endpoint cannot be missed.

\item{T-SAF-012: No personal information sent to outside services\\}

Purpose: children's and adults' personal information must not leak to
outside services, including the refinement service.

Type: Dynamic, manual. Covers SAF-SHA-002.

Initial State: an account whose adult name and email, and the child
profile's identifier, are unique marker strings.

Input: the tester uses every feature, including refinement, while recording
all network traffic.

Output: no request to an outside service contains any marker string, and
requests go only to approved destinations.

Test Case Derivation: the fit criterion of SAF-SHA-002 (H-SHA-002). Unique
markers make any leak easy to find.

How test will be performed: the test runs in Chrome on a laptop; the
DevTools Network panel is exported as a HAR file and searched for the
markers and for destinations that are not approved.

\item{T-SAF-013: Logout after inactivity\\}

Purpose: an adult who walks away from a shared device must not leave their
account and their children's data open.

Type: Dynamic, automatic. Covers SAF-SHA-006.

Initial State: an adult is logged in.

Input: (a) no activity for \texttt{IDLE\_TIMEOUT}; (b) activity just before
\texttt{IDLE\_TIMEOUT}, then no activity.

Output: the adult is logged out after \texttt{IDLE\_TIMEOUT} of
inactivity and not before, and any later request using that session is
refused.

Test Case Derivation: SAF-SHA-006 (H-SHA-001). Case (b) checks that
activity restarts the timer.

How test will be performed: Playwright uses a controlled clock to move time
forward and checks the session.

\item{T-SAF-014: Feedback does not draw the missing part\\}

Purpose: feedback should show the child what is missing without drawing it
for them, so the record shows only the child's own work.

Type: Dynamic, automatic. Covers SAF-MOU-001.

Initial State: a guided lesson at each predefined difficulty.

Input: \texttt{square-\allowbreak{}missing-\allowbreak{}side.json}, an
attempt with one side missing.

Output: the feedback shows where the side is missing, and the strokes on
screen and in the record are the same before and after the feedback.

Test Case Derivation: the fit criterion of SAF-MOU-001 (H-MOU-004).

How test will be performed: Playwright compares the stroke data before and
after the feedback is shown.

\item{T-SAF-015: Assisted, independent and creative work kept apart\\}

Purpose: the adult must be able to tell what the child drew with an
on-screen guide from what they drew alone, so progress is not overstated.

Type: Dynamic and static, manual. Covers SAF-MOU-002.

Initial State: a prepared history of one lesson record holding 10 assisted
attempts, 10 independent attempts, and 5 creative space drawings made
separately.

Input: every summary and record view is opened.

Output: assisted and independent attempts are counted and shown separately,
no creative space drawing appears, and no view ranks or rates the child by
the number of attempts.

Example: a lesson record with 10 successful attempts while tracing and 2
while drawing alone is summarized as ``10 with a guide, 2 without'', never as
``12 successful drawings''.

Test Case Derivation: the fit criterion of SAF-MOU-002 (H-MOU-005). The
assisted or independent label belongs to each attempt, not to the lesson:
following FR-LRN-012, an attempt is assisted when it is drawn with an
on-screen guide (the tracing or endpoints stage) and independent when it is
drawn with no guide. One lesson record normally holds both kinds. Help from
an adult outside Squigglr, such as guiding the child's hand, is not detected
or recorded, and this test does not cover it.

How test will be performed: a tester checks each view against the prepared
history, and a Vitest test of the summary function compares the counts.

\item{T-SAF-016: No more than three flashes per second\\}

Purpose: rapid flashing can trigger seizures, so no child-facing screen may
flash more than three times in any one second.

Type: Static (code review) and dynamic, manual. Covers SAF-MOU-003.

Initial State: the code for every animation and transition on child-facing
screens, and Squigglr open on the iPad with a guided lesson ready.

Input: (a) a team member who did not write it reviews every animation and
transition used on child-facing screens, which are all defined in
Squigglr's own code and style sheets; (b) on the iPad, the screens that
animate (all feedback, the lesson-end message, the inactivity check-in and
the screen transitions) are screen-recorded.

Output: (a) no animation or transition changes brightness or colour more
than three times in any one second; (b) stepping through each recording
frame by frame, no one-second window has more than three flashes.

Test Case Derivation: the fit criterion of SAF-MOU-003 and WCAG 2.2
success criterion 2.3.1 \citep{wcag22}. Because every animation comes from
our own code, reviewing that code finds a flashing effect before it reaches
a screen; the recordings then confirm the result on the device, without
recording every screen.

How test will be performed: once per revision (Rev~0 and Rev~1), as part of
the safety inspection; any problem is recorded as a GitHub issue labelled
\texttt{safety}.

\item{T-SAF-017: Check-in and lesson ending after inactivity\\}

Purpose: a child may walk away during a lesson, and on a shared device
another child could carry on under their profile, so an idle lesson must
check in with the child and then end, with a clear record for the adult.

Type: Dynamic, automatic and manual. Covers SAF-RAY-009.

Initial State: a guided lesson at the endpoints stage, with
\texttt{LESSON\_IDLE\_PROMPT} and \texttt{LESSON\_IDLE\_END} at the lowest
values in their ranges; the test is then repeated with the highest values.

Input: (a) no input for \texttt{LESSON\_IDLE\_PROMPT}, then a touch; (b) no
input for \texttt{LESSON\_IDLE\_END}; (c) two failed attempts, then no input
until the lesson ends.

Output: (a) the check-in appears at \texttt{LESSON\_IDLE\_PROMPT}, the touch
dismisses it, and the lesson continues at the same stage; (b) the check-in
appears, then the lesson ends at \texttt{LESSON\_IDLE\_END}, the progress is
saved, and the adult's view shows the lesson ended for inactivity; (c) the
record shows exactly two failed attempts, so the idle time added none.

Example: with illustrative values of 60 seconds for the check-in and 120
seconds for the end, a child who stops drawing sees an animated character at 60 seconds; if nobody
touches the screen, the lesson ends at 120 seconds and the adult sees ``ended:
no activity''.

Test Case Derivation: the fit criterion of SAF-RAY-009 (H-RAY-006). Case (c)
checks that inactivity is never counted as failure.

How test will be performed: Playwright uses a controlled clock to move time
forward and checks the screen and the saved record; a tester checks by hand
on the iPad that the check-in can be understood without reading and does
not flash.

\end{enumerate}

\subsection{Traceability Between Test Cases and Requirements}

\wss{Provide a table that shows which test cases are supporting which
  requirements.}

\section{Unit Test Description}

\wss{This section should not be filled in until after the MIS (detailed design
  document) has been completed.}

\wss{Reference your MIS (detailed design document) and explain your overall
philosophy for test case selection.}  

\wss{To save space and time, it may be an option to provide less detail in this section.  
For the unit tests you can potentially layout your testing strategy here.  That is, you 
can explain how tests will be selected for each module.  For instance, your test building 
approach could be test cases for each access program, including one test for normal behaviour 
and as many tests as needed for edge cases.  Rather than create the details of the input 
and output here, you could point to the unit testing code.  For this to work, you code 
needs to be well-documented, with meaningful names for all of the tests.}

\subsection{Unit Testing Scope}

\wss{What modules are outside of the scope.  If there are modules that are
  developed by someone else, then you would say here if you aren't planning on
  verifying them.  There may also be modules that are part of your software, but
  have a lower priority for verification than others.  If this is the case,
  explain your rationale for the ranking of module importance.}

\subsection{Tests for Functional Requirements}

\wss{Most of the verification will be through automated unit testing.  If
  appropriate specific modules can be verified by a non-testing based
  technique.  That can also be documented in this section.}

\subsubsection{Module 1}

\wss{Include a blurb here to explain why the subsections below cover the module.
  References to the MIS would be good.  You will want tests from a black box
  perspective and from a white box perspective.  Explain to the reader how the
  tests were selected.}

\begin{enumerate}

\item{test-id1\\}

Type: \wss{Functional, Dynamic, Manual, Automatic, Static etc. Most will
  be automatic}
					
Initial State: 
					
Input: 
					
Output: \wss{The expected result for the given inputs}

Test Case Derivation: \wss{Justify the expected value given in the Output field}

How test will be performed: 
					
\item{test-id2\\}

Type: \wss{Functional, Dynamic, Manual, Automatic, Static etc. Most will
  be automatic}
					
Initial State: 
					
Input: 
					
Output: \wss{The expected result for the given inputs}

Test Case Derivation: \wss{Justify the expected value given in the Output field}

How test will be performed: 

\item{...\\}
    
\end{enumerate}

\subsubsection{Module 2}

...

\subsection{Tests for Nonfunctional Requirements}

\wss{If there is a module that needs to be independently assessed for
  performance, those test cases can go here.  In some projects, planning for
  nonfunctional tests of units will not be that relevant.}

\wss{These tests may involve collecting performance data from previously
  mentioned functional tests.}

\subsubsection{Module ?}
		
\begin{enumerate}

\item{test-id1\\}

Type: \wss{Functional, Dynamic, Manual, Automatic, Static etc. Most will
  be automatic}
					
Initial State: 
					
Input/Condition: 
					
Output/Result: 
					
How test will be performed: 
					
\item{test-id2\\}

Type: Functional, Dynamic, Manual, Static etc.
					
Initial State: 
					
Input: 
					
Output: 
					
How test will be performed: 

\end{enumerate}

\subsubsection{Module ?}

...

\subsection{Traceability Between Test Cases and Modules}

\wss{Provide evidence that all modules have been considered.}
				
\bibliographystyle{plainnat}

\bibliography{../../refs/References}

\newpage

\section{Appendix}

This is where you can place additional information.

\subsection{Symbolic Parameters}

The definition of the test cases will call for SYMBOLIC\_CONSTANTS.
Their values are defined in this section for easy maintenance.

\subsection{Usability Survey Questions?}

\wss{This is a section that would be appropriate for some projects.}

\newpage{}

\section{Appendix --- Generative AI Use Disclosure}

\wss{ChatGPT (version 5) was used to brainstorm ideas for the template for this disclose section.}

\subsection{AI Tools Used}

\wss{Give the name of the AI tool(s) used and the version.}

\subsection{How and Where Used}

\wss{Explain what the tool(s) were used for.  Examples include: brainstorming,
explaining a technical concept, reviewing the document, generating or modifying
text, generating or modifying diagrams, identifying errors or omissions or
inconsistencies.}

\wss{For each of the uses, identify which parts of the document were affected.}

\wss{You do not need to report every interaction, but you should disclose the 
uses that materially contributed to your work}

\subsection{Verification of AI-Generated Content}

\wss{\begin{itemize}
    \item Describe how the team checked AI-generated or AI-assisted material
    before including it in the document.
    \item In particular, verify factual claims, technical assertions,
    requirements, calculations, references, and other information for which an
    AI system may produce incorrect or fabricated results.
    \item \textbf{The team is responsible for the accuracy, completeness,
consistency, and appropriateness of the submitted document, regardless of
whether material was generated by a human or an AI system.}
\end{itemize}}

\newpage{}

\section{Appendix --- Reflection}

\wss{This section is not required for CAS 741}

The information in this section will be used to evaluate the team members on the
graduate attribute of Lifelong Learning.

\input{../Reflection.text}

\begin{enumerate}
  \item What went well while writing this deliverable? 
  \item What pain points did you experience during this deliverable, and how
    did you resolve them?
  \item What knowledge and skills will the team collectively need to acquire to
  successfully complete the verification and validation of your project?
  Examples of possible knowledge and skills include dynamic testing knowledge,
  static testing knowledge, specific tool usage, Valgrind etc.  You should look to
  identify at least one item for each team member.
  \item For each of the knowledge areas and skills identified in the previous
  question, what are at least two approaches to acquiring the knowledge or
  mastering the skill?  Of the identified approaches, which will each team
  member pursue, and why did they make this choice?
\end{enumerate}

\end{document}